\section{Experiment 51: design across redshift}
\label{sec:exp51_methods}

\ExpFiftyOne{} extends the direct conversion to five epochs and asks whether it
can be made portable. Its tests are ordered from mechanical representation
checks to increasingly strong predictive claims. The main sample contains 2395
galaxies; the epoch-local and concentration analyses use the 2365-galaxy common
sample. All model comparisons use five held-out galaxy folds and 32 stochastic
draws unless stated otherwise.

\subsection{Why the CoG coordinates changed}

A small development subset did not contain the most compact valid object in the
full five-epoch sample. During the full-sample representation preflight, one
$z=2$ galaxy was found to contain 81.1 per cent of its measured stellar mass
inside 2 kpc. Its ordinary $R_{20}$ and $R_{50}$ are both below the first radial
grid point. Assigning $R_{20}=R_{50}=2\kpc$ would violate strict ordering;
discarding the object would remove exactly the type of compact galaxy the model
must retain.

The outer-mass definition in \cref{eq:outer_level,eq:exp51_coordinates} solves
this measurement problem without adding a parametric inner profile. The first
coordinate still fixes the 2-kpc fraction, while the other three radii describe
the resolved distribution of the remaining mass. The representation is
revalidated independently at every epoch before fitting any halo relation.

\subsection{Independent single-epoch ceilings}

At each redshift, the readable linear and degree-2 models are fitted
independently and compared with a three-component PCA model on identical folds.
The readable representation passes an epoch if
\begin{equation}
  \mathrm{CRPS}_{\rm readable}(z)-\mathrm{CRPS}_{\rm PCA}(z)
  \leq s_{\rm fold,PCA}(z).
  \label{eq:exp51_parity_gate}
\end{equation}
This stage answers whether a direct map is feasible at each snapshot. It does
not yet assert that one relation evolves smoothly with redshift.

At $z=2$, additional controls compare final descendant mass alone, concurrent
halo mass alone, complete descendant \diffmah{}, shuffled descendant MAH shape,
and complete descendant \diffmah{} plus concentration. These controls reveal
whether a superficially rich final-descendant feature vector is actually more
informative than the halo mass at the correct epoch.

\subsection{Descendant-conditioned and epoch-local maps}
\label{sec:two_scopes}

The two prediction scopes are shown in \cref{fig:information_boundary}. For a
population selected at $z=0.4$, the descendant-conditioned relation is
\begin{equation}
  P\!\left[\vect{\theta}_{\cog}(z)\given
  \vect{\theta}_{\rm MAH}^{\,z=0.4},\ctwo(z=0.4)\right].
  \label{eq:descendant_map}
\end{equation}
The complete history is a property of the selected descendant, even when the
target is its earlier progenitor. This model is useful for reconstructing
histories of a low-redshift-selected sample.

For an independently selected catalog at redshift $z$, the portable relation is
\begin{equation}
  P\!\left[\vect{\theta}_{\cog}(z)\given
  \vect{\theta}_{\rm MAH}^{\,\leq t_z;t_z},\ctwo(z)\right],
  \label{eq:local_map}
\end{equation}
where \diffmah{} is fitted only to $t\leq t_z$ and anchored at $t_z$. The
notation $\leq t_z;t_z$ states both the data boundary and the fit endpoint.

\begin{figure}[htbp]
  \centering
  \begin{tikzpicture}[
      font=\small,
      epoch/.style={circle,draw,minimum size=7mm,inner sep=0pt},
      box/.style={draw,rounded corners,align=center,minimum height=11mm,minimum width=34mm},
      arrow/.style={-{Latex[length=2.2mm]},thick}
    ]
    \node[font=\bfseries,hsblue] at (0,2.2) {Descendant-conditioned};
    \draw[thick,hsblue] (-3,1.2) -- (3,1.2);
    \foreach \x/\z in {-2.5/2.0,-1.25/1.5,0/1.0,1.25/0.7,2.5/0.4}{
      \node[epoch,fill=white] (d\x) at (\x,1.2) {};
      \node[below=2mm] at (\x,1.2) {$z=\z$};
    }
    \draw[very thick,hsblue,-{Latex[length=2.5mm]}] (-2.5,1.2) -- (2.5,1.2);
    \node[box,draw=hsblue,below=13mm] (dfit) at (0,1.2) {one complete-history fit\\anchored at $z=0.4$};
    \draw[arrow,hsblue] (2.5,1.2) -- (dfit.north east);

    \node[font=\bfseries,hsgreen] at (0,-2.0) {Epoch-local};
    \draw[thick,hsgreen] (-3,-3.0) -- (3,-3.0);
    \foreach \x/\z in {-2.5/2.0,-1.25/1.5,0/1.0,1.25/0.7,2.5/0.4}{
      \node[epoch,fill=white] at (\x,-3.0) {};
      \node[below=2mm] at (\x,-3.0) {$z=\z$};
      \draw[very thick,hsgreen,-{Latex[length=1.8mm]}] (-3,-2.72) -- (\x,-2.72);
    }
    \node[align=center,font=\footnotesize,text=hsgray] at (0,-4.05)
      {each prediction uses a new fit truncated and anchored at that epoch};
  \end{tikzpicture}
  \caption{Information boundary in \ExpFiftyOne{}. The descendant-conditioned
  model may use later growth because the population is selected by its final
  descendants. The epoch-local model is restricted to information available by
  each prediction epoch and is the portable form of the conversion.}
  \label{fig:information_boundary}
\end{figure}

The official halo history is truncated at every epoch and fitted for every
galaxy with $t\geq1\gyr$. This produces $2365\times5=11825$ epoch-local
\diffmah{} fits. A current-halo-mass-only baseline and mass-conditioned MAH-shape
and concentration shuffles are repeated at each snapshot.

\subsection{Held-epoch coefficient closure}

Smooth-looking coefficient curves are not sufficient evidence for a
continuous-redshift model. For each interior epoch, that epoch is omitted, the
descendant-conditioned relation is fitted at the other four snapshots, and its
coefficients are interpolated quadratically to the missing redshift. The
interpolated prediction is compared with both a direct fit and the nearest
available snapshot. The predeclared criterion is
\begin{equation}
  \frac{\mathrm{CRPS}_{\rm interp}(z)}
       {\mathrm{CRPS}_{\rm direct}(z)}-1 < 0.05.
  \label{eq:closure_gate}
\end{equation}
Because the descendant features are identical at all target epochs, one
training-fold feature standardization is meaningful across this test. The same
procedure cannot yet be transferred mechanically to the epoch-local features,
whose distributions and physical anchors evolve with redshift.

\subsection{Concentration at the correct epoch}

With descendant \diffmah{} held fixed, three concentration descriptions are
compared: $\ctwo(z=0.4)$, concurrent $\ctwo(z)$, and all earlier concentration
measurements available by $z$. Concurrent and historical concentration values
are separately shuffled within concurrent halo-mass bins. This identifies
information in the halo--galaxy pairing rather than improvement caused solely
by extra columns or their marginal distributions.

\subsection{Temporal residual coherence}
\label{sec:ar1}

Independent snapshot draws would produce discontinuous individual histories,
while reusing one residual vector at all epochs would make histories too
persistent. The out-of-fold coordinate residuals are standardized and their
cross-epoch correlation is summarized by
\begin{equation}
  \operatorname{Corr}(\epsilon_{a,i},\epsilon_{a,j})
  \approx \rho^{|i-j|},
  \label{eq:ar1}
\end{equation}
where $a$ indexes a profile coordinate and $i,j$ index the ordered snapshots.
The fitted scalar $\rho$ is used to draw an AR(1)-in-epoch latent sequence, while
the within-epoch coordinate covariance retains the model in
\cref{eq:covariance}. Coherence is judged using the Spearman rank correlation
of total mass, $R_{50}$, and $f_2$ between snapshots. This tests galaxy histories
directly; no collection of single-epoch scores can replace it.

\subsection{Symbolic transfer at \texorpdfstring{$z=2$}{z=2}}

The final stage repeats the nested-fold residual symbolic search using only the
epoch-local $z=2$ inputs. Four models are compared: linear, the three stable
$z=0.4$ terms, a new $z=2$ symbolic search, and the dense degree-2 relation.
The purpose is not to maximize the symbolic search budget. It is to test whether
the low-complexity nonlinear structure discovered at $z=0.4$ transfers to a
different epoch and input definition.
